\Chapter{Singspiele}
%\addthumb{Singspiele}{\space\Huge\sffamily\bfseries \thechapter}{white}{bbe}
\newpage
%%%%%%%%%%%%%%%%%%%%%%%%%%%%%%%%%%%%%%%%%%%%%%%%%%%%%%%%%%%%%%%%%%%%%%%%%%%%%%%%%%%%%%%%%%%%%%%%%%%%
%Pferderennen
%%%%%%%%%%%%%%%%%%%%%%%%%%%%%%%%%%%%%%%%%%%%%%%%%%%%%%%%%%%%%%%%%%%%%%%%%%%%%%%%%%%%%%%%%%%%%%%%%%%%
\Section{ Pferderennen }
\begin{table}[H]
\begin{tabular}{>{\bfseries}lp{8cm}}
   Gruppengröße:   & 5-99 Spieler*innen\\
    Alter: & 6 \\
\end{tabular}
\end{table} 
Alle Pferde gehen an den Start.„Auf die Plätze, fertig, los (Startschuss)!“ →Mit den Händen Pistole zeigen \newline
und Schießgeräusch machenDie Pferde rennen. → Klatschen auf die Oberschenkel \newline
Sie springen über ein Hindernis. → Beim Sprung über
das Hindernis Hände hochreißen Sie rennen weiter.\newline 
Da kommt eine Kurve. → Mit dem Körper in die Kurve legenSie galoppieren weiter. \newline
Am Rand winken Zuschauer. → WinkenDie Pferde rennen weiter.\newline
Da sind Fotografen  →Fotoapparat darstellenund klick, klick, klick machen\newline
Plötzlich reiten sie durch Matsch → Mitdem Mund pft, pft, pft machen\newline
Die Pferde sind auch der Zielgeraden und laufen nochmals ganz schnell. \newline
Noch 10 Meter bis zum Ziel.Noch 9 Meter. → abwärts zählen 10, 9, 8, 7, 6, ...Geschafft !!! \newline
Das Pferd mit der Startnummer „9“ ist der Sieger.\newline
Der Sieger reitet eine Ehrenrunde und winkt den Zuschauern\\



%%%%%%%%%%%%%%%%%%%%%%%%%%%%%%%%%%%%%%%%%%%%%%%%%%%%%%%%%%%%%%%%%%%%%%%%%%%%%%%%%%%%%%%%%%%%%%%%%%%%
%Dackel Waldemar
%%%%%%%%%%%%%%%%%%%%%%%%%%%%%%%%%%%%%%%%%%%%%%%%%%%%%%%%%%%%%%%%%%%%%%%%%%%%%%%%%%%%%%%%%%%%%%%%%%%%
\Section{ Dackel Waldemar }
\begin{table}[H]
\begin{tabular}{>{\bfseries}lp{8cm}}
   Gruppengröße:   & 10-99 Spieler*innen\\
    Alter: & 10 \\
\end{tabular}
\end{table} 

Dackel Waldemar ist ein Singspiel. Der Text wird gesungen oder auch gesprochen und es werden passende Bewegungen dazu gemacht.Mein Dackel Waldemar und ich, wir zwei, (auf sich zeigen, mit der Hand fast auf den Boden gehen, zwei mit zwei Fingern anzeigen), wohnen (mit den Händen ein Dach über dem Kopf machen) in der Regenbogenstraße (einen Regenbogen mit den Händen andeuten)3 (mit den Fingern die Zahl 3 zeigen), und wenn wir Abends eine Runde gehen (auf der Stelle laufen), dann kann man Dackelbeine wackeln sehen,. (wie Dackelbeine mit den Beinen wackeln). \\



%%%%%%%%%%%%%%%%%%%%%%%%%%%%%%%%%%%%%%%%%%%%%%%%%%%%%%%%%%%%%%%%%%%%%%%%%%%%%%%%%%%%%%%%%%%%%%%%%%%%
%Jack und Tina
%%%%%%%%%%%%%%%%%%%%%%%%%%%%%%%%%%%%%%%%%%%%%%%%%%%%%%%%%%%%%%%%%%%%%%%%%%%%%%%%%%%%%%%%%%%%%%%%%%%%
\Section{ Jack und Tina }
\begin{table}[H]
\begin{tabular}{>{\bfseries}lp{8cm}}
   Gruppengröße:   & 10-99 Spieler*innen\\
    Alter: & 6 \\
\end{tabular}
\end{table} 
Jack saß in der Küche mit Tina, \newline
Jack saß in der Küche mit Tina,\newline
Jack saß in der Küche mit Tina,\newline
und sie spielten auf dem alten Banjo.\newline
Sie spielten: \newline
flie, fla, fiedeleiho, \newline
flie, fla, fiedeleiho,\newline
flie, fla, fiedeleiho,\newline
und sie spielten auf dem alten Banjo.\newline\newline

Dabei werden typische Bewegungen ausgeführt:\newline
\newline
Bei "Jack saß in der Küche mit Tina" wird im Takt der Musik abwechselnd auf die eigenen Oberschenkel und auf die der Nachbarn geschlagen (zuerst eigene, dann die vom rechten Nachbarn, dann eigene, dann die vom linken Nachbarn).Bei und sie spielten auf dem alten Banjo: Banjo/Gitarre-Spiel nachahmen flie: eine Querflöte wird rechts vom Mund gespielt fla: eine Querflöte links vom Mund gespielt fiedeleiho: eine Geige wird gespielt

Üblicherweise werden mehrere Runden mit steigender Geschwindigkeit gespielt.\\



%%%%%%%%%%%%%%%%%%%%%%%%%%%%%%%%%%%%%%%%%%%%%%%%%%%%%%%%%%%%%%%%%%%%%%%%%%%%%%%%%%%%%%%%%%%%%%%%%%%%
%Funky Chicken
%%%%%%%%%%%%%%%%%%%%%%%%%%%%%%%%%%%%%%%%%%%%%%%%%%%%%%%%%%%%%%%%%%%%%%%%%%%%%%%%%%%%%%%%%%%%%%%%%%%%
\Section{ Funky Chicken }
\begin{table}[H]
\begin{tabular}{>{\bfseries}lp{8cm}}
   Gruppengröße:   & 5-99 Spieler*innen\\
    Alter: & 6 \\
\end{tabular}
\end{table} 
Alle Spieler stehen in einem Kreis. Einer beginnt den folgenden Dialog, die Gruppe antwortet: „Let me see a ...“ –„What did you say?“ –„Let me see a ...“ –„What did you say?“ –„I said [Bewegung]“ Der Spieler, der angefangen hat, macht das Tier/Wesen, das er vorher gerufen hat nach und alle anderen machen mit. Die Bewegung wird dreimal wiederholt. Dann gehen alle zurück in den Kreis. Und der gleiche oder auch jemand anders fängt den Dialog wieder an. Es können unter anderemfolgende Tiere gespielt werden(beim Ausgestaltender Figuren ist eure Kreativität gefragt):
    \begin{itemize}
        \item Funky Chicken“: Hünchenflügel mit Armen und Geräusch: „Boag-Boag
        \item Crazy Monkai: Kratzen und Geräusch: Ug-Ug
        \item Toilet Diver: Nase zuhalten und mit der anderen Hand an der Toilettenspülung ziehen Geräusch: Zwusch
        \item wildtiger: Krallen mit den Händen, Geräusch: Fauch
    \end{itemize}


%%%%%%%%%%%%%%%%%%%%%%%%%%%%%%%%%%%%%%%%%%%%%%%%%%%%%%%%%%%%%%%%%%%%%%%%%%%%%%%%%%%%%%%%%%%%%%%%%%%%
%Fli Fly Flow
%%%%%%%%%%%%%%%%%%%%%%%%%%%%%%%%%%%%%%%%%%%%%%%%%%%%%%%%%%%%%%%%%%%%%%%%%%%%%%%%%%%%%%%%%%%%%%%%%%%%
\Section{ Fli Fly Flow }
\begin{table}[H]
\begin{tabular}{>{\bfseries}lp{8cm}}
   Gruppengröße:   & 5-99 Spieler*innen\\
    Alter: & 6 \\
\end{tabular}
\end{table} 

    Fli /: \newline
Fli Fly/: \newline
Fli fly flow /:\newline
Wiste /:\newline
Gumma nana gumma nana gumma nana wiste /:\newline
oh no no no no no wiste/:\newline
exa mini sana mini u a a guana mini /:\newline
exa mini sana mini u a a gua/:\newline
hey bei de Hotten Totten forn ma aus de linken Socken /:\newline
fornma aus de Schua /:\newline

Der Text wird jeweils vom Vorsänger gerufen und anschließend von der Gruppe wiederholt. Das gesamte Stück wird meistmehrfach mit ansteigender Geschwindigkeit wiederholt.
    \\



%%%%%%%%%%%%%%%%%%%%%%%%%%%%%%%%%%%%%%%%%%%%%%%%%%%%%%%%%%%%%%%%%%%%%%%%%%%%%%%%%%%%%%%%%%%%%%%%%%%%
%A ram sam sam
%%%%%%%%%%%%%%%%%%%%%%%%%%%%%%%%%%%%%%%%%%%%%%%%%%%%%%%%%%%%%%%%%%%%%%%%%%%%%%%%%%%%%%%%%%%%%%%%%%%%
\Section{ A ram sam sam }
\begin{table}[H]
\begin{tabular}{>{\bfseries}lp{8cm}}
   Gruppengröße:   & 5-99 Spieler*innen\\
    Alter: & 6 \\
\end{tabular}
\end{table} 


    A ram sam sam, \newline
 a ram sam sam \newline
Guli guli guli guli guli ram sam sam\newline
A ram sam sam,\newline
 a ram sam sam\newline
Guli guli guli guli guli ram sam sam\newline
A ra-vi, a ra-vi\newline
Guli guli guli guli guli ram sam sam\newline
A ra-vi, a ra-vi\newline
Guli guli guli guli guli ram sam sam\newline\newline

Bewegungen\newline

Es gibt zahllose Varianten der Bewegungen, die ausgeführt werden können. Hier wird nur eine Möglichkeit beispielhaft angegeben.
A ram sam sam -Fäuste bilden und zusammenschlagen; zuerst rechts über links, dann umgekehrt
Guli guli -Flache Hände vor die Brust strecken und dort Kreisen (Winken)
A ra-vi -Mit dem Zeigefinger Richtung Kopf zeigen und drehen, zum Abschluss nach oben in den Himmel deuten.
    \\




%%%%%%%%%%%%%%%%%%%%%%%%%%%%%%%%%%%%%%%%%%%%%%%%%%%%%%%%%%%%%%%%%%%%%%%%%%%%%%%%%%%%%%%%%%%%%%%%%%%%
%The big fat pony
%%%%%%%%%%%%%%%%%%%%%%%%%%%%%%%%%%%%%%%%%%%%%%%%%%%%%%%%%%%%%%%%%%%%%%%%%%%%%%%%%%%%%%%%%%%%%%%%%%%%
\Section{ The big fat pony }

\begin{table}[H]
\begin{tabular}{>{\bfseries}lp{8cm}}
   Gruppengröße:   & 5-99 Spieler*innen\\
    Alter: & 6 \\ 
\end{tabular}
\end{table} 
 Die Spieler stehen im Kreis, nicht zu eng, so dass sich jeder frei bewegen kann. Es wird der Liedertext ("here we go with the big fat pony") gesungen, alle bereits "aktiven" Spieler (zuerst nur der erste Spieler) laufen dabei im Kreis. Anschließend bleiben die aktiven Spieler bei (jeweils) einem noch passiven Spieler stehen und führen die Bewegungen mit diesem gemeinsam durch. Diese bisherpassiven Spieler werden dadurch zu aktiven und laufen in der nächsten Runde ebenfalls im Kreis mit.
Liedertext und BewegungenDie Melodie des Liedes entspricht dem bekannten Lied What shall we do with the drunken sailor?:
Here we go with the big fat pony,here we go with the big fat pony,here we go with the bigfat pony,Die aktiv laufenden Spieler laufen im Uhrzeigersinn im Kreis 
early in the morning.und stellen sich zu einem noch passiven SpielerFront, front, front, my baby,Die jeweils zusammenstehenden aktiven und passiven Spieler shaken (rhythmisches Disco-Tanzen), wobei sie Brust an Brust stehen,back, back, back, my baby,... sie Hintern an Hintern stehen ...side, side, side, my baby, ... sie Seite an Seite stehen.early in the morning. \\



%%%%%%%%%%%%%%%%%%%%%%%%%%%%%%%%%%%%%%%%%%%%%%%%%%%%%%%%%%%%%%%%%%%%%%%%%%%%%%%%%%%%%%%%%%%%%%%%%%%%
%If you happy and youknow it 
%%%%%%%%%%%%%%%%%%%%%%%%%%%%%%%%%%%%%%%%%%%%%%%%%%%%%%%%%%%%%%%%%%%%%%%%%%%%%%%%%%%%%%%%%%%%%%%%%%%%
\Section{ If you happy and youknow it  }

\begin{table}[H]
\begin{tabular}{>{\bfseries}lp{8cm}}
   Gruppengröße:   & 5-99 Spieler*innen\\
    Alter: & 6 \\ 
\end{tabular}
\end{table} 


    In der großen Gruppe wird das Lied gesungen. Die Haupt-Phrase ("clap your hands") wird dabei mit einer Bewegung unterstützt, hier also in die Hände geklatscht. \newline 

If you're happy and you know it, clap your hands.\newline
If you're happy and you know it, clap your hands.\newline
If you're happy and you know it,\newline
And you really want to show it,\newline
If you're happy and you know it, clap your hands.\newline
In der nächsten Strophe wirddie Phrase geändert, also etwa durch stomp your feet ausgetauscht;\newline entsprechend wird dann nicht geklatscht, sondern zweimal mit den Füßen aufgestampft.\newline

weitere Strophensnap your fingers (zweimal mit den Fingern schnippen)\newline
shout "Hurray!" (Schreien "hoo-ray")turn around (einmal um die eigene Achse drehen)\newline
slap your knees (beide Knie zusammen schlagen)\newline
do it all (alle zuvor gemachten Bewegungen hintereinander)\newline
    \\



%%%%%%%%%%%%%%%%%%%%%%%%%%%%%%%%%%%%%%%%%%%%%%%%%%%%%%%%%%%%%%%%%%%%%%%%%%%%%%%%%%%%%%%%%%%%%%%%%%%%
%Singingin the rain
%%%%%%%%%%%%%%%%%%%%%%%%%%%%%%%%%%%%%%%%%%%%%%%%%%%%%%%%%%%%%%%%%%%%%%%%%%%%%%%%%%%%%%%%%%%%%%%%%%%%
\Section{ Singingin the rain }
\begin{table}[H]
\begin{tabular}{>{\bfseries}lp{8cm}}
   Gruppengröße:   & 5-99 Spieler*innen\\
    Alter: & 6 \\ 
\end{tabular}
\end{table} 
I'm singing in the rain \newline
just singing in the rain\newline
what a glorious feeling\newline
I'm happy again\newline
\newline
Zuerst wird die Strophe aus dem Lied gesungen. Anschließend ruft der Betreuer (er unterbricht dadurch das Lied, das ja noch weiter ginge):\newline

Hands up   (deutsch: "Arme nach oben")\newline

Die Gruppe wiederholt die Anweisung und führt diese ebenfalls aus.\newline

Anschließend wird am Stand locker getanzt (mit den Hüften gewackelt) und dazu folgender Text gesprochen/gesungen:\newline\newline

A zumzaza, A zumzaza, A zumzazaza-aha\newline
A zumzaza, A zumzaza, Azumzazaza-aha\newline

Danach beginnt die nächste Runde. Es wird wieder die Liederstrophe gesungen, das Lied vom Betreuer unterbrochen. Dann wiederholt er alle bisherigen Bewegungen und fügt eine weitere hinzu. In der ersten Runde führt die Gruppe also eine besondere Bewegung vor dem Tanz aus, in der zweiten Runde zwei und so weiter. Besonders lustigwird dies, wenn etwa die Zunge aus dem Mund gestreckt werden muss.\newline

Bewegungen\newline\newline
Folgende Bewegungen werden vor dem Tanz ausgeführt, und müssen während des Tanzes beibehalten werden!
\newline
Hands up (Hände nach vorne strecken)\newline
Thumbs up (Daumen nach oben, "OK"-Geste)\newline
Elbows back (Ellenbogen zurück, das heißt an die Hüfte anlegen, während die Hände nach vorne gestreckt bleiben)\newline
Knees together (Knie zusammendrücken)\newline
Head back (Kopf nach hinten)\newline
Tongue out (Zunge rausstrecken)\newline
        \\

%%%%%%%%%%%%%%%%%%%%%%%%%%%%%%%%%%%%%%%%%%%%%%%%%%%%%%%%%%%%%%%%%%%%%%%%%%%%%%%%%%%%%%%%%%%%%%%%%%%%
%Karawanen -Song (Sum gali gali)
%%%%%%%%%%%%%%%%%%%%%%%%%%%%%%%%%%%%%%%%%%%%%%%%%%%%%%%%%%%%%%%%%%%%%%%%%%%%%%%%%%%%%%%%%%%%%%%%%%%%
\Section{ Karawanen -Song (Sum gali gali) }

\begin{table}[H]
\begin{tabular}{>{\bfseries}lp{8cm}}
   Gruppengröße:   & 5-99 Spieler*innen\\
    Alter: & 6 \\
\end{tabular}
\end{table}
    Die Gruppe wird in drei Untergruppen eingeteilt. Die eine Gruppe sind die Kamele die machen: Uhh –ahhpuhh(dabei gehen die Arme erst hoch, dann nach außen gestreckt und dann pustest man Luft durch die Lippen aus und senkt seinen Kopf) Die andere Gruppe sind die singenden Mädchen, sie singen: laalalalalaaaalalaalalalalaaala(dazu machen sie Bauchtanz Bewegungen mit den Armen). Die letzte Gruppe sind die Kameltreiber, die machen: Sum gali gali gali sum galigali Jetzt steuert man mit unterschiedlichen Lautstärken und Geschwindigkeiten die Gruppe und bekommt so ein witziges Lied hin. Besonders gut funktioniert es auch, wenn man die einzelnen Gruppen alleine singen lässt.
    \\

\Section{Auf dem Donnerbalken saßen zwei Gestalten}
\begin{table}[H]
\begin{tabular}{>{\bfseries}lp{8cm}}
	Gruppengröße:   & 5-99 Spieler*innen\\
	Alter: & 6 \\
\end{tabular}
\end{table}

Auf dem Donnerbalken saßen zwei Gestalten \\
und sie schrien nach Klopapier, Klopapier. \\

Und dann kam der zweite, der sich zu ihm reihte\\
und sie schrien nach Klopapier, Klopapier\\

Und da kam der Dritte, setzt sich in die Mitte \\
und sie schrien nach Klopapier, Klopapier. \\

Und da kam der Vierte, als die Scheiße schmierte \\ 
und sie schrien nach Klopapier, Klopapier. \\

Und da kam der Fünfte, der die Nase rümpfte\\
und sie schrien nach Klopapier, Klopapier.\\

Und da kam der Sechste, als die Scheiße kleckste\\ 
sie schrien nach Klopapier, Klopapier.\\

Und da kam der Siebte, als der Balken wippte\\
und sie schrien nach Klopapier, Klopapier.\\

Und da kam der Achte, als der Balken krachte\\
und sie schrien nach Klopapier, Klopapier.\\

Und da kam der Neunte, als die Scheiße schäumte\\
und sie schrien nach Klopapier, Klopapier.\\

Und da kam der Zehnte, brachte das ersehnte\\
KLO-PA-PIER!\\

(Und dann kam der elfte, nam sich gleich der Hälfte,\\
und sie schrien nach Klopapier! Klopapier! Klopapier!)


